% GENERATED FILE. Edit canonical lessons, then run npm run generate:books.
% canonical-source-hash: fnv1a64:e1af1ed322778927
% canonical-lessons: GU-C330-mandap, GU-C330-ghumvu, GU-C330-jhulvu, GU-C330-bhajan, GU-C330-bhakt

\chapter{Nine Nights of Dance}
\label{ch:gu-nine-nights-of-dance}
\hlchaptermodality{330}

\begin{chapteropening}
\textbf{By the end of this chapter:} \emph{I can tell how people danced and sang at night.}

Complete the last lesson of chapter 330: I can tell how people danced and sang at night.
\end{chapteropening}

\section[\texorpdfstring{\gu{મંડપ} (manḍap)}{મંડપ (manḍap)}]{\gu{મંડપ} (manḍap) — a pavilion (for a festival)}

\label{lesson:GU-C330-mandap}



\begin{quote}
Before the new one: say the Gujarati for a garland, then the Gujarati for a procession.
\end{quote}

\subsection*{You'll want to know: \gu{મંડપ}}
\textbf{\gu{મંડપ}} — \emph{manḍap} — \textquotedblleft{}a pavilion (for a festival)\textquotedblright{}.

\begin{quote}
\gu{ચોકમાં} \gu{મોટો} \gu{મંડપ} \gu{હતો}.
\end{quote}

\begin{quote}
\emph{cokmā̃ moṭo manḍap hato.}
\end{quote}

\begin{quote}
There was a big pavilion in the square.
\end{quote}

A masculine noun.

\begin{cousinweb}
\textbf{\gu{મંડપ}} (\emph{manḍap}) is a tatsama, Sanskrit \emph{maṇḍapa}, \textquotedblleft{}an open temporary shed, a pavilion\textquotedblright{} (CDIAL 9740). It is put up for weddings and festivals.
\end{cousinweb}

\begin{grammarlens}[title={What you've built}]
One more word for telling what happened.
\end{grammarlens}

\subsection*{Your turn}
\begin{itemize}
\raggedright
  \item \emph{Say it:} \emph{manḍap}
  \item \emph{Say it:} \emph{manḍap}, once more
  \item \emph{Say it:} \emph{cokmā̃ moṭo manḍap hato.} — There was a big pavilion in the square.
  \item \emph{Recall:} say the Gujarati for a garland, then the Gujarati for a procession, then say \emph{manḍap} again
\end{itemize}

\begin{tcolorbox}[breakable,colback=teal!4,colframe=teal!35!black,title={Before you move on}]
What does \gu{મંડપ} mean? (\textquotedblleft{}A pavilion (for a festival)\textquotedblright{}.) Say it once more.
\end{tcolorbox}
\section[\texorpdfstring{\gu{ઘૂમવું} (ghūmvũ)}{ઘૂમવું (ghūmvũ)}]{\gu{ઘૂમવું} (ghūmvũ) — to whirl round}

\label{lesson:GU-C330-ghumvu}



\begin{quote}
Before the new one: say the Gujarati for a procession, then the Gujarati for a pavilion.
\end{quote}

\subsection*{You'll want to know: \gu{ઘૂમવું}}
\textbf{\gu{ઘૂમવું}} — \emph{ghūmvũ} — \textquotedblleft{}to whirl round\textquotedblright{}.

\begin{quote}
\gu{રાત્રે} \gu{લોકો} \gu{મંડપમાં} \gu{ગરબે} \gu{ઘૂમ્યા}.
\end{quote}

\begin{quote}
\emph{rātre loko manḍapmā̃ garbe ghūmyā.}
\end{quote}

\begin{quote}
At night people whirled round in the garba in the pavilion.
\end{quote}

Past (he …): \gu{તે} \gu{ગરબે} \gu{ઘૂમ્યો}, \emph{te garbe ghūmyo}.

\begin{cousinweb}
\textbf{\gu{ઘૂમવું}} (\emph{ghūmvũ}) comes from \emph{*ghummati}, a form CDIAL reconstructs, where it lists Gujarati \emph{ghumvũ}, glossed \textquotedblleft{}to revolve in the mind\textquotedblright{} (CDIAL 4485). In the dance it is moving round and round: \textbf{\gu{ગરબે} \gu{ઘૂમ્યા}}, they danced the garba.
\end{cousinweb}

\begin{grammarlens}[title={What you've built}]
One more word for telling what happened.
\end{grammarlens}

\subsection*{Your turn}
\begin{itemize}
\raggedright
  \item \emph{Say it:} \emph{ghūmvũ}
  \item \emph{Say it:} \emph{ghūmvũ}, once more
  \item \emph{Say it:} \emph{rātre loko manḍapmā̃ garbe ghūmyā.} — At night people whirled round in the garba in the pavilion.
  \item \emph{Recall:} say the Gujarati for a procession, then the Gujarati for a pavilion, then say \emph{ghūmvũ} again
\end{itemize}

\begin{tcolorbox}[breakable,colback=teal!4,colframe=teal!35!black,title={Before you move on}]
What does \gu{ઘૂમવું} mean? (\textquotedblleft{}To whirl round\textquotedblright{}.) Say it once more.
\end{tcolorbox}
\section[\texorpdfstring{\gu{ઝૂલવું} (jhūlvũ)}{ઝૂલવું (jhūlvũ)}]{\gu{ઝૂલવું} (jhūlvũ) — to swing}

\label{lesson:GU-C330-jhulvu}



\begin{quote}
Before the new one: say the Gujarati for a pavilion, then the Gujarati for to whirl round.
\end{quote}

\subsection*{You'll want to know: \gu{ઝૂલવું}}
\textbf{\gu{ઝૂલવું}} — \emph{jhūlvũ} — \textquotedblleft{}to swing\textquotedblright{}.

\begin{quote}
\gu{બાળકો} \gu{હિંચકા} \gu{પર} \gu{ઝૂલતાં} \gu{હતાં}.
\end{quote}

\begin{quote}
\emph{bāḷko hinckā par jhūltā̃ hatā̃.}
\end{quote}

\begin{quote}
The children were swinging on the swing.
\end{quote}

Past (she …): \gu{તે} \gu{છોકરી} \gu{ઝૂલી}, \emph{te chokrī jhūlī}.

\begin{cousinweb}
\textbf{\gu{ઝૂલવું}} (\emph{jhūlvũ}) comes from \emph{*jhulyati}, a form CDIAL reconstructs, and CDIAL lists the Gujarati verb, \textquotedblleft{}to swing, hang\textquotedblright{} (CDIAL 5406). People swing on a \textbf{\gu{હિંચકો}}; leaves swing on a branch.
\end{cousinweb}

\begin{grammarlens}[title={What you've built}]
One more word for telling what happened.
\end{grammarlens}

\subsection*{Your turn}
\begin{itemize}
\raggedright
  \item \emph{Say it:} \emph{jhūlvũ}
  \item \emph{Say it:} \emph{jhūlvũ}, once more
  \item \emph{Say it:} \emph{bāḷko hinckā par jhūltā̃ hatā̃.} — The children were swinging on the swing.
  \item \emph{Recall:} say the Gujarati for a pavilion, then the Gujarati for to whirl round, then say \emph{jhūlvũ} again
\end{itemize}

\begin{tcolorbox}[breakable,colback=teal!4,colframe=teal!35!black,title={Before you move on}]
What does \gu{ઝૂલવું} mean? (\textquotedblleft{}To swing\textquotedblright{}.) Say it once more.
\end{tcolorbox}
\section[\texorpdfstring{\gu{ભજન} (bhajan)}{ભજન (bhajan)}]{\gu{ભજન} (bhajan) — a hymn}

\label{lesson:GU-C330-bhajan}



\begin{quote}
Before the new one: say the Gujarati for to whirl round, then the Gujarati for to swing.
\end{quote}

\subsection*{You'll want to know: \gu{ભજન}}
\textbf{\gu{ભજન}} — \emph{bhajan} — \textquotedblleft{}a hymn\textquotedblright{}.

\begin{quote}
\gu{રાત્રે} \gu{મારા} \gu{દાદાએ} \gu{ભજન} \gu{ગાયાં}.
\end{quote}

\begin{quote}
\emph{rātre mārā dādāe bhajan gāyā̃.}
\end{quote}

\begin{quote}
At night my grandfather sang hymns.
\end{quote}

A neuter noun.

\begin{cousinweb}
\textbf{\gu{ભજન}} (\emph{bhajan}) is a tatsama, Sanskrit \emph{bhajana}, \textquotedblleft{}worship, adoration\textquotedblright{} (Monier-Williams): a song of devotion. It is neuter: \textbf{\gu{ભજન} \gu{ગાયાં}}, sang hymns.
\end{cousinweb}

\begin{grammarlens}[title={What you've built}]
One more word for telling what happened.
\end{grammarlens}

\subsection*{Your turn}
\begin{itemize}
\raggedright
  \item \emph{Say it:} \emph{bhajan}
  \item \emph{Say it:} \emph{bhajan}, once more
  \item \emph{Say it:} \emph{rātre mārā dādāe bhajan gāyā̃.} — At night my grandfather sang hymns.
  \item \emph{Recall:} say the Gujarati for to whirl round, then the Gujarati for to swing, then say \emph{bhajan} again
\end{itemize}

\begin{tcolorbox}[breakable,colback=teal!4,colframe=teal!35!black,title={Before you move on}]
What does \gu{ભજન} mean? (\textquotedblleft{}A hymn\textquotedblright{}.) Say it once more.
\end{tcolorbox}
\section[\texorpdfstring{\gu{ભક્ત} (bhakt)}{ભક્ત (bhakt)}]{\gu{ભક્ત} (bhakt) — a devotee}

\label{lesson:GU-C330-bhakt}



\begin{quote}
Before the new one: say the Gujarati for to swing, then the Gujarati for a hymn.
\end{quote}

\subsection*{You'll want to know: \gu{ભક્ત}}
\textbf{\gu{ભક્ત}} — \emph{bhakt} — \textquotedblleft{}a devotee\textquotedblright{}.

\begin{quote}
\gu{મંદિરમાં} \gu{ઘણા} \gu{ભક્ત} \gu{આવ્યા}.
\end{quote}

\begin{quote}
\emph{mandirmā̃ ghaṇā bhakt āvyā.}
\end{quote}

\begin{quote}
Many devotees came to the temple.
\end{quote}

A masculine or feminine noun, by the person meant.

\begin{cousinweb}
\textbf{\gu{ભક્ત}} (\emph{bhakt}) is a tatsama, Sanskrit \emph{bhakta}, \textquotedblleft{}devoted to, attached to\textquotedblright{} (Monier-Williams). The same Sanskrit word, in the sense of a share of food, gave \textbf{\gu{ભાત}}, rice.
\end{cousinweb}

\begin{grammarlens}[title={What you've built}]
One more word for telling what happened.
\end{grammarlens}

\subsection*{Your turn}
\begin{itemize}
\raggedright
  \item \emph{Say it:} \emph{bhakt}
  \item \emph{Say it:} \emph{bhakt}, once more
  \item \emph{Say it:} \emph{mandirmā̃ ghaṇā bhakt āvyā.} — Many devotees came to the temple.
  \item \emph{Recall:} say the Gujarati for to swing, then the Gujarati for a hymn, then say \emph{bhakt} again
\end{itemize}

\begin{tcolorbox}[breakable,colback=teal!4,colframe=teal!35!black,title={Before you move on}]
What does \gu{ભક્ત} mean? (\textquotedblleft{}A devotee\textquotedblright{}.) Say it once more.
\end{tcolorbox}
